% !TEX root = ../main-es.tex
%TC:macro \limite [ignore,ignore]
%TC:macro \nota [ignore]
%TC:macro \TODO [ignore]
%TC:macro \parencite [ignore]
%TC:macro \textcite [ignore]
\section{Introducción}\label{sec:introduccion}

\subsection{Motivación}

Existe hoy una clase emergente de sistemas a la que la literatura ya le puso nombre, los \emph{agentes auto-evolutivos} \parencite{gao2026selfevolving}, en la que el agente modifica durante su propia ejecución los componentes que lo definen, ya sea el flujo de trabajo que sigue, la memoria que consulta o el programa mismo que lo está ejecutando, de forma que la auto-modificación dejó de ser una especulación teórica para volverse una propiedad de sistemas desplegados. Las instancias de esa clase son públicas y verificables: Voyager programa sus propias habilidades en Minecraft y las repara cuando su verificación falla \parencite{wang2023voyager}, Reflexion reescribe su memoria verbal tras cada fracaso \parencite{shinn2023reflexion}, Agent Workflow Memory induce flujos de trabajo reutilizables mientras opera \parencite{wang2024awm}, y AFlow y ADAS buscan automáticamente, representados como código, el flujo que el agente ejecuta \parencite{zhang2024aflow,hu2024adas}. El caso límite lo documentaron los propios autores de The AI Scientist, cuyo agente editó su código para extenderse el límite de tiempo que lo acotaba y en otra corrida se relanzó a sí mismo \parencite{lu2024aiscientist}, lo cual muestra que la capacidad de reescribirse alcanza ya a las restricciones puestas para contenerla.

El problema es que esa capacidad se despliega sin un criterio medido de cuándo ejercerla, pues los orquestadores actuales fijan por intuición cada cuánto conviene replanificar, y la pregunta por el cuándo sigue abierta en la agenda de la propia clase, que la nombra explícitamente como la dimensión \emph{when} \parencite{gao2026selfevolving}. Este trabajo aborda esa pregunta desde una instancia de calibración concreta, un motor de ejecución agéntica propio (OASys) en el que un agente basado en un modelo de lenguaje puede leer y reescribir el programa que lo ejecuta, y en el que las tres decisiones que importan se resolvieron por intuición: Las ediciones surten efecto en la ejecución siguiente (de modo que el trabajo en curso se pierde), porque nadie sabía si compensaba construir la reparación en caliente; la deliberación se acota con un tope fijo de cinco iteraciones, sin justificación derivada; y el único freno contra un ciclo de reescrituras improductivas es la autorización humana. Vale la pena anotarlo con honestidad: Se trata de un sistema propio y no publicado, de forma que aquí aporta la pregunta y los parámetros que un despliegue real podría medir, pero no evidencia, la cual recae por completo sobre el experimento reproducible que se describe en las secciones siguientes.

Al leer la teoría clásica de agentes con esa instancia en mente resulta que el mecanismo completo ya está descrito con otro vocabulario \parencite{russellNorvig2020}: El programa editable es un plan (jerárquico, §11.4), revisarlo contra lo que se observa es monitoreo de ejecución y reescribirlo es replanificación (§11.5), lo que el agente tiene a la vista es observabilidad parcial (§2.3.2), y aplicar la edición en la ejecución siguiente o sobre la ejecución en curso es la distinción entre planificar en lazo abierto y en lazo cerrado (§11.5). Si el problema es clásico, cabe la pregunta clásica: El beneficio de replanificar crece con el dinamismo del entorno, mientras que su costo (el tiempo durante el cual el agente delibera y el mundo sigue cambiando) es aproximadamente constante, por lo que debería existir un punto de cruce, y es ese punto el que nadie había medido.

\subsection{La pregunta}

\begin{quote}
  \itshape\proyPregunta
\end{quote}

Para responderla se construyó un modelo basado en agentes en MESA \parencite{kazil2020mesa}, en el que un agente recolector reúne recursos que se desplazan sobre una retícula con obstáculos, y en el que la decisión de modelado central es tratar la deliberación como una acción más, que ocupa ticks del reloj durante los cuales el mundo se sigue moviendo. El replanificador no es un modelo de lenguaje real sino un oráculo abstracto con dos parámetros, la competencia $p$ (la probabilidad de que la reparación propuesta sea la óptima) y el costo $c$ (los ticks que tarda en proponerla), cuyos errores son plausibles y no ruido. Cabe mencionar que esa abstracción es una decisión metodológica y no solo de presupuesto, pues convierte una propiedad del modelo de lenguaje en un parámetro del experimento: Con un modelo real solo se puede responder si la auto-edición funcionó con ese modelo en ese momento, mientras que con el oráculo se puede preguntar a partir de qué competencia y de qué costo conviene dejar que un agente se reescriba, que es un enunciado que sobrevive al cambio de modelo.

\subsection{Aportes}

\begin{enumerate}[leftmargin=1.6em, itemsep=0.3em]
  \item \textbf{Un modelo reproducible de la auto-edición de planes,} con ocho arquitecturas que recorren la taxonomía de agentes de \textcite{russellNorvig2020}, desde el reflejo simple hasta dos agentes que razonan sobre su propia creencia antes de reparar, con la deliberación como acción interrumpible y con una regla de costo que cobra el mismo replanificador al mismo precio a toda política que planifica. El modelo se acompaña de 278 pruebas y de once controles de sanidad, la mayoría de ellos exactos corrida a corrida.
  \item \textbf{Un barrido de 16.260 corridas} en trece barridos, con 30 semillas pareadas por combinación, más 600 series temporales y 3.000 simulaciones de calibración, cuyo hallazgo central es que la frontera que importa no separa al lazo cerrado del lazo abierto (que casi empatan) sino a deliberar de no deliberar.
  \item \textbf{La cuantificación del valor de reparar en caliente} frente a reiniciar en frío, que es la decisión de ingeniería que motivó el trabajo, junto con su robustez al orden de activación, su evolución en el tiempo y su dependencia de la abundancia de objetivos.
  \item \textbf{Dos resultados negativos sobre saber qué se sabe:} Un agente basado en conocimiento con fluentes, cuya ventaja proviene de reparar menos y se invierte en cuanto preguntar cuesta un tick, y un agente bayesiano cuya creencia está calibrada pero cuya decisión es peor que la del agente lógico, porque su regla de utilidad esperada repara de más.
  \item \textbf{Un criterio operable de cuándo reescribir,} que consiste en estimar el dinamismo $\lambda$ a partir de la telemetría de un despliegue por computación bayesiana aproximada y leer la recomendación en la frontera ya barrida, con la advertencia de que la competencia $p$ del replanificador no es identificable con esa telemetría.
  \item \textbf{Un registro honesto de una corrección metodológica,} la de la regla de costo, que mostró que una asimetría de contabilidad entre políticas había fabricado varias de las conclusiones de una versión anterior del modelo.
\end{enumerate}

\subsection{Estructura del documento}

La sección~\ref{sec:marco} reúne el marco conceptual (agentes y entornos de tarea, planificación en lazo abierto y cerrado, costo de la deliberación, agentes basados en conocimiento e inferencia probabilística), mientras que la sección~\ref{sec:modelo} describe el modelo, incluidas la regla de costo, el oráculo y la formalización de los dos agentes que infieren. La sección~\ref{sec:diseno} fija el diseño experimental, con los barridos, los controles de sanidad y las hipótesis escritas antes de ver los datos, y la sección~\ref{sec:resultados} presenta los resultados con sus errores estándar y su veredicto. Progresivamente, la sección~\ref{sec:avance2} contrasta el resultado final con lo que anticipaba el Avance~2, la sección~\ref{sec:discusion} discute por qué la frontera tiene la forma que tiene y qué criterio de decisión se desprende, la sección~\ref{sec:limites} recoge los límites y, finalmente, la sección~\ref{sec:conclusiones} cierra con las conclusiones.
